%% kapa: exam=1 round=30 year=2019 period=1 form=A subject=civil_law subject_ko=민법 seq=17 number=17
%% source: Q-Net 감정평가사 기출문제 (cst00309) artlSeq=5205205 file=extracted/1008/5205205/2019년 제30회 감정평가사 1차시험 문제지/2019년 제30회 감정평가사 1차시험 문제지/2019년도 제30회 감정평가사 1차 1교시 A형.hwp
%% conversion: kordoc 4.24.51 -> tools/gen1.py (auto); review=proofread-sonnet; flags=none
\begin{question}{17}
조건과 기한에 관한 설명으로 옳지 않은 것은?
\begin{choices}
  \item 조건성취의 효력은 당사자의 의사표시로 소급하게 할 수 없다.
  \item 조건이 법률행위 당시에 이미 성취할 수 없는 것일 경우에는 그 조건이 해제조건이면 조건없는 법률행위로 한다.
  \item 조건의 성취가 미정인 권리는 일반규정에 의하여 담보로 할 수 있다.
  \item 당사자의 특약이나 법률행위의 성질상 분명하지 않으면 기한은 채무자의 이익을 위한 것으로 추정한다.
  \item 기한의 이익은 포기할 수 있지만, 이로 인해 상대방의 이익을 해하지 못한다.
\end{choices}
\end{question}
\answer{1}
